\section{Introduction}
\label{sec:intro}

Large-scale service platforms, such as telecom and cloud operators or IoT back-ends, have two
problems the small version of the same system does not have. First, \emph{finding things}: with
hundreds of instances spread over sites, an alert is useless unless it says which instance, on
which machine, in which data centre. Second, \emph{defending at scale}: the defences must work on
every instance at once, must not overload the links between sites with monitoring traffic, and must
not hurt the system they protect.

The objective of this project is to design and run a service-based architecture that scales to
large scale, monitor it with Prometheus, and integrate it with a SIEM (Wazuh) that detects three
types of incident from the course syllabus and responds automatically. The system is a small IoT
pipeline in which temperature sensors send readings to an ingestion service, which stores them
through a storage service, while an analytics service computes statistics and a registry provides
naming and discovery. The platform runs as Docker images in a two-node Kubernetes (k3s) cluster in
which each node is a different data centre. An earlier deployment with one Incus container per
service was kept as a baseline, and every test runs unchanged on both.

\paragraph{Contributions.}
\begin{itemize}
  \item A hierarchical naming scheme (\emph{instance.node.rack.dc.region.country}) that is used
        unchanged by the registry, the metrics, the SIEM agents and the logs, so the source of any
        event can be located by name (Section~\ref{sec:arch}).
  \item Location-aware discovery and per-location Prometheus aggregates, which reduce the
        federated monitoring traffic to 0.7\% of a raw federation (Section~\ref{sec:eval}).
  \item Detection and automatic response for a massive HTTP denial of service, junk email and DNS
        hijacking, including a Kubernetes-specific block that acts before \code{kube-proxy}
        (Section~\ref{sec:threats}).
  \item Correlation of anomaly events into a single root-cause alert, following the FCAPS fault
        management workflow (Section~\ref{sec:risk}).
  \item Automated proof-of-concept tests that record detection and response times, false
        positives and recovery as evidence, and an honest account of what failed.
\end{itemize}

\paragraph{Overview.}
Section~\ref{sec:sota} reviews related work and tools. Section~\ref{sec:arch} describes the
architecture and naming. Section~\ref{sec:risk} analyses the attack surface, presents an
attack--defense tree and the monitoring strategy. Section~\ref{sec:threats} covers the three
threats, Section~\ref{sec:eval} evaluates the approaches and their scalability, and
Section~\ref{sec:concl} concludes. The appendix lists how to reproduce every result.
